We evaluate the two-stage procedure on simulated Gaussian data against four benchmarks: Stage~1 alone (decomposable empirical Bayes), the birth-death MCMC sampler of \textcite{mohammadi2015bdgraph}, the graphical lasso of \textcite{friedman2008sparse}, and its semiparametric extension, the nonparanormal estimator of \textcite{liu2009nonparanormal}. The birth-death sampler, hereafter BDMCMC, is the central comparator.\footnote{BDMCMC is the default sampler of the \texttt{BDgraph} R package \parencite{mohammadi2015bdgraph}; the configuration we use throughout this paper (a uniform Bernoulli prior on the graph with $g_{ij} = 0.2$, a G-Wishart prior on the precision matrix $\Omega$, and the closed-form normalizing-constant approximation of \textcite{mohammadi2023accelerating}) corresponds to the package defaults invoked by the call \texttt{BDgraph::bdgraph(data = Y)}. The BDMCMC default is the configuration an applied researcher running the procedure as released by the authors actually faces.} It is the leading inferential tool for Bayesian Gaussian graphical models over the unrestricted graph space, and it is the procedure the DIP-augmented two-stage estimator extends.

% BDMCMC represents the joint posterior over $(G, \Omega)$ as a continuous-time Markov jump process. At each state, every candidate edge $e \in E \cup \bar{E}$ has a Poisson birth/death rate determined by the posterior ratio of flipping that edge. The next event is sampled from the discrete distribution induced by those rates and is accepted by construction, so every iteration produces a state change. Earlier implementations evaluated the rate ratios via an inner Monte Carlo or exchange-algorithm step. \textcite{mohammadi2023accelerating} replace that step with a closed-form approximation of the G-Wishart normalizing constant, reducing the per-edge cost to a single gamma-function evaluation and bringing the sampler into the problem-size range relevant for applied work. The simulation design varies graph size, topology, and Monte Carlo replicate.\\


The DIP and the BDMCMC default can be thought of as a prior family indexed by $\lambda$,
\[
    p_{ij}(\lambda, r_0) \;=\; \lambda\, \text{PIP}_{ij} \;+\; (1-\lambda)\, r_0.
\]
The family contains BDMCMC as its boundary at $(\lambda, r_0) = (0, 0.2)$: every entry of the prior matrix then equals $0.2$, and the procedure reduces to the package default. We fix $r_0 = 0.2$ throughout this section so that $\lambda$ alone indexes the deviation from this boundary, and any DIP-vs.-default gap in the results is cleanly attributable to Stage-1 information. The two-stage estimator we report uses $\lambda = 0.5$.

\subsection{Design}

For each replicate $r \in \{1, \dots, 5\}$ and each combination of size $p \in \{5, 10, 20\}$, topology (one of five), and diagonal strength $d \in \{3, 5, 10, 15\}$, we construct a true conditional independence graph $G^* = (V, E^*)$, build a precision matrix $\Omega^*$ with diagonal $d$ and off-diagonal $1.9$ on the edges of $G^*$, and draw $n = 500$ observations $y_i \sim \mathcal{N}_p(0, \Omega^{*-1})$.\footnote{If $\Omega^*$ is not positive definite (relevant for high-degree Barab\'asi--Albert hubs), we add a small multiple of the identity to bring its smallest eigenvalue to $0.1$.} The diagonal $d$ controls the signal-to-noise ratio: the implied maximum absolute partial correlation is $\rho_{\max} = 1.9 / d$, ranging from $\rho_{\max} \approx 0.63$ at $d = 3$ down to $\rho_{\max} \approx 0.13$ at $d = 15$. 
% The simulation generator's package-example default is $d = 15$.
We treat $d = 5$ ($\rho_{\max} \approx 0.38$) as the baseline calibration because it places most comparators above the EBIC graphical lasso's selection threshold while still leaving room for the Bayesian methods to differentiate; the star topology at $p \geq 10$ is a known exception we discuss in the Results subsection.
%We report the full grid below in the signal-strength subsection.
The five topologies fall into two groups. The decomposable group consists of the chain ($1\!-\!2\!-\!\cdots\!-\!p$) and the star (vertex $1$ connected to all others). The non-decomposable group adds the cycle ($1\!-\!2\!-\!\cdots\!-\!p\!-\!1$, which contains a chordless cycle of length $p \geq 4$); a sparse Erd\H{o}s--R\'enyi random graph with edge probability $2/p$; and a Barab\'asi--Albert preferential-attachment graph with one edge added per new node. The simulation has $300$ cases.

We fit five estimators on each case: Stage~1 of Section~\ref{sec:misspec}, the two-stage estimator at $(\lambda, r_0) = (0.5, 0.2)$, the BDMCMC default, the graphical lasso with EBIC tuning at $\gamma = 0.5$,\footnote{We tune the graphical lasso's regularization parameter via the extended Bayesian Information Criterion of \textcite{foygelEBIC2010} at $\gamma = 0.5$, the default in the \texttt{huge} R package and the EBIC strength most commonly reported in graphical-model studies. EBIC is consistent for selection but penalises edges relatively heavily, which produces the empty-graph collapse at weak signal that Section~\ref{sec:sim} (signal-strength axis) documents at $d = 15$. The Stability Approach to Regularization Selection \parencite{liu2010stability} tunes $\lambda$ by selecting the most subsampling-stable solution, tends to be more permissive than EBIC on weak signal (and more lenient on false positives), and would likely move GLASSO above the empty graph at $d = 15$ at the cost of a higher false-positive rate. A direct STARS comparison is left for follow-up work.} and the nonparanormal extension of the graphical lasso. The DGP above is Gaussian by construction, so on these simulations the nonparanormal estimator should behave like the graphical lasso: its empirical normal-scores transformation is approximately the identity on Gaussian marginals, and it reduces to running the graphical lasso on rank-transformed inputs.\footnote{We include the nonparanormal in the simulation grid for two reasons. First, as a sanity check that the semiparametric flexibility does not degrade recovery when the Gaussian assumption holds. Second, as a baseline for Section~\ref{sec:application}, where the marginals are log range volatilities and any GLASSO--nonparanormal divergence is informative about the binding force of the Gaussianity assumption.} 
%Stage 2 of the two-stage procedure and the BDMCMC default share the same MCMC budget. To isolate the method effect from sampling variability, we fit all five estimators on the same simulated dataset within each cell.

\subsection{Evaluation metric}

Each estimator returns a binary $p \times p$ adjacency matrix $\hat G$. We summarize the agreement between $\hat G$ and the truth $G^*$ using the Matthews correlation coefficient (MCC),
\begin{equation}
    \text{MCC}(\hat G, G^*) \;=\;
        \frac{TP \cdot TN - FP \cdot FN}
             {\sqrt{(TP+FP)(TP+FN)(TN+FP)(TN+FN)}},
\end{equation}
where the four counts are read off the $2 \times 2$ confusion matrix of the binary edge classification on the upper triangle of the adjacency matrix. % The MCC is the Pearson correlation between the predicted and true edge indicators, takes values in $[-1, 1]$, and equals zero for any constant predictor including the empty graph. We prefer it to the F1 score and to raw accuracy because it uses all four cells of the confusion matrix and remains informative under the heavy class imbalance of sparse graph estimation: at $p = 20$ chain there are $19$ true edges among $190$ possible, so a predictor that always returns the empty graph attains $90$ percent accuracy but $\text{MCC} = 0$.

% \subsection{Theoretical predictions}

% Theorem~\ref{thm:concentration} characterizes the asymptotic Stage-1 PIPs: $\text{PIP}_{ij} \to 1$ on $E^*$, $\text{PIP}_{ij} \to 0$ on pairs that lie in no minimal triangulation of $G^*$, and $\text{PIP}_{ij} \to$ some positive number on fill edges of competing minimal triangulations. If we focus on the DIP-vs.-default performance difference, the Stage-1 information is positive on true edges and on clear non-edges, and negative on fill edges. The net sign therefore depends on the size of the fill set $\bigcup_k F_k$ of $G^*$, which in turn depends on $|\mathcal{M}(G^*)|$.

% Three predictions follow. First, on decomposable $G^*$ ($|\mathcal{M}(G^*)| = 1$, no fill edges), the DIP transfers strictly correct information and should outperform the BDMCMC default. This covers chain, star, and the tree-like Barab\'asi--Albert ($m = 1$). Second, on non-decomposable $G^*$ with small $|\mathcal{M}(G^*)|$ (e.g., short chordless cycles, sparse near-tree structure) the DIP should still help. Erd\H{o}s--R\'enyi at $p_{\rm edge} = 2/p$ produces such graphs typically. Third, on graphs with many minimal triangulations, the chord-edge inflation dominates, and the DIP should underperform the BDMCMC default. The cycle topology at $p \geq 5$, where the number of triangulations grows combinatorially in $p$, is the cleanest test of this prediction.

\subsection{Results}

Theorem~\ref{thm:concentration} yields three predictions for the DIP-vs.-default comparison. On decomposable $G^*$ (chain, star, tree-like Barab\'asi--Albert), the Stage-1 PIPs concentrate on $E^*$ and the DIP transfers strictly correct information. On non-decomposable $G^*$ with few minimal triangulations (sparse Erd\H{o}s--R\'enyi at $2/p$ edge probability), the chord-edge inflation is bounded and the DIP should still help. On graphs with many minimal triangulations --- the cycle at $p \geq 5$, where the count grows combinatorially in $p$ --- the chord-edge inflation dominates and the DIP should underperform the BDMCMC default.

Table~\ref{tab:sim-mcc} reports mean MCC at every $(\text{topology}, \text{method}, d)$ cell, averaged over the three sizes $p \in \{5, 10, 20\}$. The baseline calibration is the $d = 5$ column ($\rho_{\max} \approx 0.38$).
\begin{table}[H]
    \small
    \centering
    \caption{Recovery performance: mean MCC by topology, method, and diagonal strength $d$. The $d = 5$ column is the baseline calibration.}
    \label{tab:sim-mcc}
    % Auto-generated by run_simulations_extended.r --- do not edit by hand.
\begin{tabular}{llrrrr}
\toprule
Topology & Method & $d=3$ & $d=5$ & $d=10$ & $d=15$ \\
 & & ($\rho_{\max}\!\approx\!0.63$) & ($\rho_{\max}\!\approx\!0.38$) & ($\rho_{\max}\!\approx\!0.19$) & ($\rho_{\max}\!\approx\!0.13$) \\
\midrule
Chain & Stage-1 & 0.89 & 0.94 & 0.79 & 0.60 \\
 & BDMCMC & 0.97 & 0.99 & 0.93 & 0.77 \\
 & Two-stage & 0.92 & 0.98 & 0.93 & 0.75 \\
 & GLASSO & 0.33 & 0.62 & 0.52 & 0.07 \\
 & NPN & 0.32 & 0.60 & 0.53 & 0.03 \\
\midrule
Star & Stage-1 & 0.99 & 0.98 & 0.83 & 0.53 \\
 & BDMCMC & 0.97 & 0.96 & 0.91 & 0.71 \\
 & Two-stage & 0.91 & 0.95 & 0.89 & 0.68 \\
 & GLASSO & 0.03 & 0.18 & 0.75 & 0.00 \\
 & NPN & 0.03 & 0.18 & 0.75 & 0.05 \\
\midrule
Cycle & Stage-1 & 0.68 & 0.57 & 0.52 & 0.58 \\
 & BDMCMC & 0.66 & 0.99 & 0.95 & 0.76 \\
 & Two-stage & 0.66 & 0.98 & 0.92 & 0.76 \\
 & GLASSO & 0.59 & 0.67 & 0.53 & 0.02 \\
 & NPN & 0.59 & 0.68 & 0.50 & 0.02 \\
\midrule
Erd\H{o}s--R\'enyi & Stage-1 & 0.86 & 0.77 & 0.54 & 0.57 \\
 & BDMCMC & 0.92 & 0.97 & 0.93 & 0.81 \\
 & Two-stage & 0.91 & 0.95 & 0.90 & 0.79 \\
 & GLASSO & 0.42 & 0.59 & 0.40 & 0.19 \\
 & NPN & 0.42 & 0.58 & 0.39 & 0.05 \\
\midrule
Barab\'asi--Albert & Stage-1 & 0.87 & 0.95 & 0.83 & 0.55 \\
 & BDMCMC & 0.97 & 0.98 & 0.96 & 0.78 \\
 & Two-stage & 0.96 & 0.98 & 0.95 & 0.78 \\
 & GLASSO & 0.19 & 0.47 & 0.54 & 0.00 \\
 & NPN & 0.19 & 0.46 & 0.56 & 0.05 \\
\bottomrule
\end{tabular}

    \\[0.4em]
    \begin{minipage}{0.92\linewidth}
        \footnotesize
        %\emph{Notes:} Mean Matthews correlation coefficient between the recovered graph $\hat G$ and the true graph $G^*$, averaged over three sizes ($p \in \{5, 10, 20\}$) and five Monte Carlo replicates per cell. Higher is better; the maximum is one. Lower $d$ corresponds to stronger signal ($\rho_{\max} = 1.9 / d$). Decomposable topologies (chain, star) are listed first, then non-decomposable topologies (cycle, Erd\H{o}s--R\'enyi, Barab\'asi--Albert). Stage-1 is the decomposable empirical Bayes posterior of \textcite{donnet2012empirical}; BDMCMC is the birth-death sampler of \textcite{mohammadi2015bdgraph} at its default uniform Bernoulli prior ($g_{ij} = 0.2$, no Stage-1 information); Two-stage is the DIP-augmented Stage-2 sampler at $(\lambda, r_0) = (0.5, 0.2)$, which nests the BDMCMC default at $\lambda = 0$.
    \end{minipage}
\end{table}

% Three patterns hold. First, both Stage-2 estimators dominate Stage-1 alone. At $d = 5$ BDMCMC averages $0.98$ across topologies and the two-stage estimator $0.97$, against $0.84$ for Stage-1 alone, a gap of about $0.14$ MCC. The decomposability restriction is the binding constraint, and lifting it through any non-decomposable Bayesian sampler recovers the bulk of the gap to ceiling MCC. Reading across $d$, the Stage-1 disadvantage widens at weak signal: Stage-1's topology-averaged MCC moves from $0.86$ at $d = 3$ to $0.57$ at $d = 15$, while BDMCMC moves from $0.90$ to $0.77$. The widening concentrates on the non-decomposable topologies where Theorem~\ref{thm:concentration}'s chord-edge inflation is largest. On cycle and Erd\H{o}s--R\'enyi, Stage-1 trails BDMCMC by $0.18$ to $0.43$ MCC across the grid. On the decomposable and tree-like topologies (chain, star, Barab\'asi--Albert) Stage-1 stays within $0.10$ MCC of BDMCMC at $d \leq 5$ but falls $0.18$ to $0.23$ behind at $d = 15$.

% Second, the calibrated DIP-vs.-default gap is small. At $d = 5$ the two-stage estimator and BDMCMC differ by at most $0.02$ MCC at any topology (mean gap $-0.01$). Reading across $d$, the gap stays within $\pm 0.06$ at every $(d, \text{topology})$ cell, with the largest absolute differences at $d = 3$ on chain ($-0.05$) and star ($-0.06$), both with BDMCMC ahead. The intuition that a structured Stage-1 prior should add value at low signal-to-noise, where the data alone leaves the Stage-2 posterior under-determined, does not survive in our designs. Stage 1 and Stage 2 read the same dataset, so when the data are too weak to constrain Stage 2 they are also too weak to constrain Stage 1's PIPs, and mixing those noisy PIPs into Stage 2 imports noise rather than information.\footnote{We probed this directly with three follow-up experiments held out from the main grid: a mixing-speed comparison at fixed $(n, p)$ to ask whether the DIP-warmed chain reaches its asymptote in fewer iterations; a sample-size sweep over $n \in \{100, 250, 500, 1000, 2000\}$ at $p = 10$ to ask whether small $n / p$ unlocks a DIP advantage; and a ``distance-from-decomposability'' axis perturbing the chain on $p = 10$ by $k \in \{0, 1, 2, 4, 8\}$ random long-distance edges. None delivers a sustained recovery advantage. The mixing-speed comparison shows a brief warm-start window for $K \lesssim 150$ Stage-2 iterations, well below typical applied budgets, after which BDMCMC overtakes; the $n$ and $k$ axes show no detectable trend in the DIP gap, consistent with the structural argument above.} The BDMCMC default itself realizes the natural ``uninformative-prior Stage~2'' alternative to the DIP; it operates over the unrestricted graph space and uses the flat $g_{ij} = 0.2$ Bernoulli prior.

% Third, the graphical lasso and the nonparanormal estimator lag the Bayesian methods at every signal level and collapse toward the empty graph at $d = 15$. At $d = 5$ GLASSO averages $0.51$ across topologies, with the lag concentrated on the star ($0.18$): the hub's degree-$(p-1)$ structure produces a highly heterogeneous partial-correlation spectrum that the EBIC penalty does not pick out as a coherent group, so the empty graph wins for nearly every $p \geq 10$ replicate. The other four topologies sit between $0.47$ (Barab\'asi--Albert) and $0.67$ (cycle) at $d = 5$, meaningful structure but well below the Bayesian methods. Reading across $d$, the lag is roughly stable through $d = 10$ and then collapses at $d = 15$: GLASSO drops to $0.07$ (chain), $0.00$ (star), $0.02$ (cycle), $0.19$ (Erd\H{o}s--R\'enyi), and $0.00$ (Barab\'asi--Albert). Four of the five values sit at or near the empty-graph baseline; only on Erd\H{o}s--R\'enyi does GLASSO retain visibly positive structure ($0.19$), and that residual is within Monte Carlo variability. The nonparanormal is essentially identical to GLASSO across the grid, as expected on the Gaussian DGP we use here (its rank-transformation reduces to the identity on Gaussian marginals). The Bayesian--frequentist gap is therefore not an artifact of the baseline calibration; at the package-example default the frequentist methods return little usable structure while the Bayesian methods retain $\text{MCC} \approx 0.71$--$0.81$. This is the regime where the informative HIW-class prior does its mechanistic work: when the likelihood alone is too weak to commit to edges, a structured prior carries the inference.
Three patterns hold. First, both Stage-2 estimators dominate Stage-1 alone: at $d = 5$, BDMCMC averages 0.98 across topologies and the two-stage estimator 0.97, against 0.84 for Stage-1 alone. The decomposability restriction is the binding constraint. Second, the DIP-vs.-default gap is negligible: at $d = 5$, the two-stage estimator and BDMCMC differ by at most 0.03 MCC at any topology. The disadvantage is most notable on the cycle ($-0.03$), but the advantage on decomposable topologies is muted because BDMCMC already obtains almost 1.00.\footnote{We probed this directly with three follow-up experiments held out from the main grid: a mixing-speed comparison at fixed $(n, p)$ to ask whether the DIP-warmed chain reaches its asymptote in fewer iterations; a sample-size sweep over $n \in \{100, 250, 500, 1000, 2000\}$ at $p = 10$ to ask whether small $n / p$ unlocks a DIP advantage; and a ``distance-from-decomposability'' axis perturbing the chain on $p = 10$ by $k \in \{0, 1, 2, 4, 8\}$ random long-distance edges. None delivers a sustained recovery advantage. The mixing-speed comparison shows a brief warm-start window for $K \lesssim 150$ Stage-2 iterations, well below typical applied budgets, after which BDMCMC overtakes; the $n$ and $k$ axes show no detectable trend in the DIP gap.} Third, the graphical lasso lags the Bayesian methods substantially (0.51 on average) and collapses on the star topology (0.18), where the hub structure produces a heterogeneous partial-correlation spectrum that EBIC does not capture as a coherent group. 
The Bayesian methods outperform the frequentist baselines by a wide margin, plausibly because the maximum partial correlation in the design is modest ($\rho_{\max} \approx 0.38$ at $d = 5$); the EBIC penalty crosses below most edges before it crosses below the empty-graph baseline.
\subsection{Discussion}

% Three takeaways. First, the simulation's purpose is to characterize the DIP family, not to find a $\lambda > 0$ that dominates the default. The family contains the BDMCMC default as the boundary $(\lambda, r_0) = (0, 0.2)$, and an auxiliary sensitivity sweep over $\lambda \in [0,1]$ on a representative subset (chain, cycle, Erd\H{o}s--R\'enyi at $p = 10$, $n = 500$; Appendix~\ref{sec:lambda}) finds $\lambda = 0$ at or near the recovery optimum on every topology in the sweep, with monotone or near-monotone degradation as $\lambda$ rises. The size of that degradation is small ($0.06$ to $0.08$ MCC at $\lambda = 1$ versus $\lambda = 0$, depending on topology), and on a truly decomposable $G^*$ with enough data to identify the triangulation cleanly the curves should flatten further. We read this as evidence that $\lambda = 0$ is the safe default for a practitioner with no positive prior on decomposability: the recovery cost of $\lambda > 0$ is small, the procedure also pays the Stage-1 compute, and the data alone via BDMCMC suffice on our designs. A practitioner who does hold a positive prior on decomposability has principled reasons to set $\lambda > 0$ regardless. Decomposable graphs are more tractable, they admit closed-form normalizing constants, and they inherit Theorem~\ref{thm:concentration}'s asymptotic guarantees, and the empirical penalty in our designs is small.

% Second, the two-stage procedure is conceptually more expensive than the BDMCMC default by construction, since it runs Stage 1 in addition to a Stage-2 call that is essentially the BDMCMC default at a non-flat prior. Stage 1 also has structural features that push the per-iteration gap further: its Metropolis-Hastings proposals can be rejected (so many iterations leave the graph unchanged) while the BDMCMC default's birth-death events are accepted by construction, and its acceptance ratio reads marginal-likelihood ratios between HIW posteriors via clique- and separator-level Cholesky operations while BDMCMC has the closed-form normalizing-constant approximation of \textcite{mohammadi2023accelerating}. Both features are intrinsic to the two estimators rather than implementation artifacts. The implication is that the DIP procedure earns its compute budget through what its two-graph output tells the practitioner, not through any improvement in recovery quality.

% Third, what its two-graph output tells the practitioner is a misspecification diagnostic for Theorem~\ref{thm:concentration}'s asymptotic guarantees. The decomposable Stage-1 graph $\hat G_D$ is a controlled estimator of a minimal triangulation of $G^*$ under the theorem's assumptions; the unrestricted Stage-2 graph $\hat G$ is the practitioner's actual point estimate. If $\hat G \approx \hat G_D$ (or $\hat G$ is itself decomposable), the asymptotic theorem cleanly applies and qualitative claims about the recovered structure inherit those guarantees. If $\hat G$ introduces many edges that close chordless cycles relative to $\hat G_D$, the practitioner is outside that regime and the cycle-closing edges should be read as conjectural rather than asymptotically grounded. The DDLY application in Section~\ref{sec:application} runs this diagnostic explicitly.

% The main caveat is that the design spans $p \in \{5, 10, 20\}$, well below the dimensions in which the graphical lasso is typically used. Readers interested in the high-dimensional regime should consult the comparative review of \textcite{mohammadi2015bdgraph}. Section~\ref{sec:application} examines the procedure at $p = 106$ on the DDLY bank-volatility data, where partial correlations are estimated from observed series rather than constructed.


The simulations' purpose is to characterize the decomposability assumption, both alone (Stage 1) and as an informative prior (two-stage). Appendix~\ref{sec:lambda} sweeps $\lambda$ over $[0,1]$ on a representative subset of topologies; the recovery curves are nearly flat, with at most $0.06$--$0.08$ MCC between $\lambda = 0$ and $\lambda = 1$.

What the two-stage procedure adds is diagnostic rather than higher accuracy. The decomposable Stage-1 graph $\hat G_D$ is a controlled estimator of a minimal triangulation of $G^*$ under the assumptions of Theorem~\ref{thm:concentration}; the unrestricted Stage-2 graph $\hat G$ is the practitioner's actual point estimate. If $\hat G \approx \hat G_D$ (or $\hat G$ is itself decomposable), the asymptotic theorem cleanly applies and qualitative claims about the recovered structure inherit those guarantees. If $\hat G$ introduces many edges that close chordless cycles relative to $\hat G_D$, the practitioner is outside that regime and the cycle-closing edges should be read as conjectural rather than asymptotically grounded. Section~\ref{sec:application} runs this diagnostic explicitly on the DDLY application.

Stage-1 alone performs better than the worst-case bound suggests, which indicates that the decomposability restriction is not particularly costly in moderate-$p$ settings. The main caveat is that our simulation spans $p \le 20$, well below the dimensions where the GLASSO is typically used. In our application we test it at $p = 106$.